% SPDX-License-Identifier: Apache-2.0
\section{An MDIO Master}
\label{sec:x-mdio}

An Ethernet PHY keeps its configuration in 32 registers of 16 bits,
and MDIO is how a MAC reaches them. It is the management interface of
IEEE 802.3, clause 22: a clock, MDC, which the master always drives,
and one data line, MDIO, which the master drives except while the PHY
answers a read. The line has a pull-up, so a line nobody drives reads
one.

A frame is 64 bits. The first 32 are ones, the preamble. Then come the
start, \code{01}; the operation, \code{10} for a read and \code{01}
for a write; the PHY's address and the register's, five bits each;
two bits of turnaround; and the sixteen bits of the word. On a write
the master drives all of it. On a read it lets go of the line at the
turnaround, and the PHY drives a zero and then the word.

\code{Mdio} has four registers, and one frame per command.

\begin{center}\footnotesize
\begin{tabular}{@{}llp{0.52\columnwidth}@{}}
\toprule
Offset & Name & What it is \\
\midrule
\code{0x00} & \code{ctrl} & the divider \\
\code{0x04} & \code{cmd} & one frame; writing it starts the frame \\
\code{0x08} & \code{data} & the word the last read took \\
\code{0x0c} & \code{state} & busy, and done \\
\bottomrule
\end{tabular}
\end{center}

Half of an MDC cycle is \code{div} cycles and one. The divider resets
to zero, which is far faster than any PHY takes, so a driver writes
\code{ctrl} first.

\subsection{When the Line Is Read}

A bit goes onto the line while MDC is low, and the PHY takes it as MDC
rises. A bit the PHY drives changes just after MDC rises, and 802.3
gives the PHY up to 300\,ns to do it. The master therefore takes each
bit at the end of the low half, just before MDC rises again, a whole
MDC cycle after the PHY changed it. At 1\,MHz that leaves a thousand
nanoseconds for the PHY's three hundred.

\subsection{How It Is Written}

The master is written the way the I2C master of
Section~\ref{sec:x-i2c} is: two processes. The bus process answers the
four registers every cycle and counts the cycles of a half in
\code{tick}. The engine waits for a write to \code{cmd} and loads two
registers of 32 bits from it. \code{frame} holds the bits after the
preamble. \code{drive} holds a one for each bit the master drives: all
of a write, and the first fourteen bits of a read. Then come two
counted loops of 32 bits, the preamble and the frame, each bit a low
half and a high half. The two registers shift together, so the line's
level and whether it is driven both come from the top bit. Incoming frame
bits sampled from the data line shift into \code{rx}; at the conclusion of a
read cycle, the final sixteen bits represent the returned register word.

\subsection{What Checks It}

\code{MdioPhy} is a model of a PHY at one address with 32 registers,
stepped between cycles from the loop. It takes a bit as MDC rises,
waits for 32 ones and a start, and answers a read to its address.
Seven tests in \code{//lib/parts:parts\_test} drive the master
against it:
\begin{itemize}
\item a read of both identifier words;
\item a read of all 32 registers;
\item an address nobody answers, which reads all ones;
\item a write, read back;
\item a register on a vendor page, read through the page select in
register 31, and the old page written back;
\item a check that the master and the PHY never drive the line in the
same cycle, and that every half of MDC lasts at least \code{div} and
one cycles;
\item the lowering.
\end{itemize}

The example scans for its PHY the way a driver does. It reads the
first identifier register at each address until one answers with
something other than all ones, then reads the second identifier and
the status register. The model holds the identifier and the status
word the JL2121 on the AX7A200B answered with, \code{937c 4032} and
\code{796d}.

\begin{figure}[htbp]
\centering
\input{mdio_timing.tex}
\caption{The end of a read frame: the master drives the address bits,
lets go of the line at the turnaround, and the PHY's word shifts into
\code{rx} a bit per MDC cycle.}
\label{fig:x-mdio}
\end{figure}

\lstinputlisting[caption={\code{ex\_mdio.rs}. Run with \code{bazel run //lib/examples:ex\_mdio}.},label={lst:x-mdio}]{lib/examples/ex_mdio.rs}

\lstinputlisting[style=ascii,caption={What \code{ex\_mdio} printed when the build ran it: an address with no PHY, then the PHY found and read.},label={lst:x-mdio-out}]{out_mdio.txt}
